\documentclass{article}

\usepackage[utf8]{inputenc}
\usepackage[spanish]{babel}
\usepackage{tikz}

% Ya no necesitamos 'backgrounds' ni 'fit'
\usetikzlibrary{shapes,arrows,positioning} 

\tikzset{
	block/.style = {
		draw, 
		fill=blue!10, 
		rectangle, 
		minimum height=3em, 
		minimum width=5em, 
		line width=1pt, 
		align=center, 
		rounded corners=2pt 
	},
	signal/.style = {
		->, 
		>=stealth', 
		thick, 
		shorten >=1pt, 
		shorten <=1pt
	},
	label/.style = {
		font=\small, 
		midway, 
		above,
		align=center
	}
}

\begin{document}
	
	\begin{figure}[htbp]
		\centering
		\begin{tikzpicture}[node distance=2.5cm, auto]
			
			\node (input) {};
			
			\node [block, right=of input] (sys1) {Sistema 1 \\ $S_1$};
			
			\node [block, right=of sys1] (sys2) {Sistema 2 \\ $S_2$};
			
			\node [right=of sys2] (output) {};
			
			% --- Conexiones ---
			\draw [signal] (input) -- node[label] {Entrada, $x(t)$} (sys1);
			\draw [signal] (sys1) -- node[label] {Señal \\ Intermedia, $w(t)$} (sys2);
			\draw [signal] (sys2) -- node[label] {Salida, $y(t)$} (output);
			
			% --- CAJA PUNTEADA CORREGIDA ---
			% Dibujamos el rectángulo tomando como referencia la esquina superior izquierda de S1
			% y la esquina inferior derecha de S2, dándoles un margen (xshift/yshift).
			\draw[dashed, gray, thick, rounded corners=5pt]
			([xshift=-0.7cm, yshift=0.8cm]sys1.north west) rectangle 
			([xshift=0.7cm, yshift=-1.5cm]sys2.south east);
			
			% Colocamos el texto exactamente en el medio de ambos bloques, desplazado hacia abajo
			\path (sys1.south east) -- (sys2.south west) 
			node[midway, below=0.6cm, font=\small] {Sistema Total Comprendido};
			
		\end{tikzpicture}
		\caption{Ejemplo de interconexión de sistemas en cascada o serie. La salida del Sistema 1 es la entrada directa del Sistema 2.}
		\label{fig:interconexion_serie}
	\end{figure}
	
\end{document}